\section{Coordinate slices in arbitrary dimension}
\label{sec:hastings05_coordinate_slices}

Hastings proposes treating a higher-dimensional system as a chain by grouping
sites according to one coordinate \cite[lines 745--756]{Hastings2006Solving}.
The construction below carries this out for finite coordinate boxes. It
constructs the pair operators from the physical local terms and gives an
explicit projection length polynomial in the logarithm of the physical number
of sites. Its coordinate setting and the corrected closing window are recorded
in \cite{gap:hastings05_coordinate_slices,gap:hastings05_hierarchical_energy}.

\begin{definition}[The operators assigned to coarse edges]
    \label{def:hastings05_edge_partition}
    \lean{TNLean.MPS.AreaLaw.LocalSolution.edgePartition,
        TNLean.MPS.AreaLaw.LocalSolution.sum_edgePartition,
        TNLean.MPS.AreaLaw.LocalSolution.norm_edgePartition_le,
        TNLean.MPS.AreaLaw.LocalSolution.isHermitian_edgePartition,
        TNLean.MPS.AreaLaw.LocalSolution.centered_edgePartition,
        TNLean.MPS.AreaLaw.LocalSolution.residual_edgePartition_le,
        TNLean.MPS.AreaLaw.LocalSolution.centeredNorm_edgePartition_le,
        TNLean.MPS.AreaLaw.LocalSolution.centeredNorm_edgePartition_le_two_mul,
        TNLean.MPS.AreaLaw.LocalSolution.edgePartition_mem_supportedOperators}
    \leanok
    Let $M_i$ be a finite family of physical terms and assign each index $i$
    to one coarse edge $f(i)$. Define $A_j=\sum_{f(i)=j}M_i$.
    Then $\sum_jA_j=\sum_iM_i$, and Hermiticity is preserved.
    If there are $n$ physical terms with $\|M_i\|\le J$, then
    $\|A_j\|\le nJ$. For a unit reference state $\psi$, write
    $c_i=\operatorname{Re}\langle\psi,M_i\psi\rangle$ and
    $c_j'=\operatorname{Re}\langle\psi,A_j\psi\rangle$.
    Local residuals at most $\epsilon$ give
    $\|(A_j-c_j')\psi\|\le n\epsilon$, and
    $\|A_j-c_j'\|\le2nJ$.
    If every term assigned to edge $j$ acts on $S_j$, so does $A_j$.
\end{definition}
\begin{proof}\leanok
    The fibres of $f$ partition the physical terms, and expectation is
    additive. Apply the triangle inequality to each fibre sum. Finally
    $|c_j'|\le\|A_j\|$, so centring costs at most a factor of two.
\end{proof}

\begin{theorem}[The geometry of coordinate slices]
    \label{thm:hastings05_coordinate_slice_geometry}
    \lean{TNLean.MPS.AreaLaw.LocalSolution.coordinateSlice,
        TNLean.MPS.AreaLaw.LocalSolution.coordinateSliceLower,
        TNLean.MPS.AreaLaw.LocalSolution.coordinateSlice_eq_lower_or_succ,
        TNLean.MPS.AreaLaw.LocalSolution.coordinateSliceEdge,
        TNLean.MPS.AreaLaw.LocalSolution.coordinateSliceEdge_support,
        TNLean.MPS.AreaLaw.LocalSolution.coordinateSlice_lt,
        TNLean.MPS.AreaLaw.LocalSolution.coordinateSliceEdge_closing_support,
        TNLean.MPS.AreaLaw.LocalSolution.natDist_coordinate_le_of_dist_le,
        TNLean.MPS.AreaLaw.LocalSolution.coordinateSliceOperator,
        TNLean.MPS.AreaLaw.LocalSolution.sum_coordinateSliceOperator,
        TNLean.MPS.AreaLaw.LocalSolution.coordinateSliceOperator_mem_supportedOperators,
        TNLean.MPS.AreaLaw.LocalSolution.coordinateSliceOperator_mem_supportedOperators_of_metric,
        TNLean.MPS.AreaLaw.LocalSolution.isHermitian_coordinateSliceOperator,
        TNLean.MPS.AreaLaw.LocalSolution.residual_coordinateSliceOperator_le,
        TNLean.MPS.AreaLaw.LocalSolution.centeredNorm_coordinateSliceOperator_le,
        TNLean.MPS.AreaLaw.LocalSolution.centeredNorm_coordinateSliceOperator_le_two_mul,
        TNLean.MPS.AreaLaw.LocalSolution.coordinateSliceOperator_closing_mem_supportedOperators}
    \leanok
    \uses{def:hastings05_edge_partition}
    Give each physical site a nonnegative integer coordinate $z(i)$ and
    assume $M_i$ acts only where $|z(v)-z(i)|\le\ell$, with $\ell\ge1$.
    Put $b(v)=\lfloor z(v)/(2\ell)\rfloor$ and
    $b_i=\lfloor\max\{0,z(i)-\ell\}/(2\ell)\rfloor$.
    The support of $M_i$ meets only slices $b_i$ and $b_i+1$.
    For $B$ coarse slices, assign $M_i$ to edge $b_i+1$ when $b_i+1<B$,
    and to edge zero otherwise. Every internal $A_j$ then acts only on
    slices $j-1,j$. If $z(v)<2\ell B$ for all sites, edge zero acts only
    on the last and first slices. The sum and estimates of the preceding
    definition hold for these actual edge operators.

    On a finite coordinate box with the coordinate maximum metric,
    metric-ball locality implies the required coordinate-strip locality.
    The implication also holds for physical metrics dominating that metric.
\end{theorem}
\begin{proof}\leanok
    The coordinate interval of a term lies between
    $\max\{0,z(i)-\ell\}$ and this lower endpoint plus $2\ell$.
    Division by $2\ell$ gives precisely the two stated slices.
    The coordinate extent excludes a slice numbered $B$, which gives the
    boundary assertion. Projection to one coordinate does not increase
    maximum-metric distance.
\end{proof}

\begin{theorem}[The complete coordinate-slice construction]
    \label{thm:hastings05_coordinate_slice_approximation}
    \lean{TNLean.MPS.AreaLaw.LocalSolution.coordinateSliceProjection,
        TNLean.MPS.AreaLaw.LocalSolution.exists_unit_mem_coordinateSliceProjection,
        TNLean.MPS.AreaLaw.LocalSolution.card_le_two_mul_re_trace_of_coordinateSliceProjection,
        TNLean.MPS.AreaLaw.LocalSolution.hierarchyProjectionLength_coordinateSlices_le}
    \leanok
    \uses{thm:hastings05_coordinate_slice_geometry,
        thm:hastings05_quantitative_periodic, thm:hastings05_retained_tree,
        thm:hastings05_explicit_projection_length}
    Let the physical site set be nonempty, of cardinality $n$, and let
    $J,\Delta>0$. In the preceding coordinate setting, let the terms be
    Hermitian with $\|M_i\|\le J$, and suppose the unit reference state
    $\psi$ satisfies $\|(M_i-c_i)\psi\|\le Je^{-r}$.
    Set $K=J_{\mathrm{eff}}=2nJ$, $h=\max\{2,\lceil\log_2q\rceil\}$,
    and take
    \begin{align*}
        r\ge512h^2\bigl(1+\max\{0,\log(2nJ/\Delta)\}\bigr).
    \end{align*}
    The constructed hierarchy, with its closing window, retains a unit
    vector $\phi$ satisfying
    \begin{align*}
        \|\phi-\psi\|\le2(3^{h+1}-1)e^{-r/8}.
    \end{align*}
    This vector belongs to every retained block space below the root,
    and the finite slice family covers every physical site.

    Suppose moreover that $H$ has ground projection $P_0$, ground energy
    $E_0$, and gap $\Delta$, that $\psi$ is a ground state of $H$, and that
    $\|H-\sum_iM_i\|\le nJe^{-r}$. Then every orthonormal family retained
    by the constructed projection has at most $2\operatorname{rank}P_0$
    members. If $B\le n$, the sufficient projection length is at most
    \begin{align*}
        512(\lceil\log_2n\rceil+2)^3
        \bigl(1+\max\{0,\log(2J/\Delta)\}\bigr).
    \end{align*}
\end{theorem}
\begin{proof}\leanok
    The edge residual is at most $nJe^{-r}\le Ke^{-r}$ and its centred
    norm is at most $K$. Apply the explicit finite hierarchy to these
    constructed edges. The closing window retains all descendant spaces.
    The physical gap gives the count after the approximation error and
    the radius sum are included. The cubic logarithmic bound follows from
    $B\le n$ and the effective coupling $J_{\mathrm{eff}}=2nJ$.
\end{proof}

\begin{corollary}[Full boxes in arbitrary dimension]
    \label{cor:hastings05_coordinate_slice_boxes}
    \lean{TNLean.MPS.AreaLaw.LocalSolution.coordinateSide_le_card_box,
        TNLean.MPS.AreaLaw.LocalSolution.coordinateBox_slice_extent,
        TNLean.MPS.AreaLaw.LocalSolution.coordinateBox_ball_subset_strip,
        TNLean.MPS.AreaLaw.LocalSolution.hierarchyProjectionLength_coordinateBox_le,
        TNLean.MPS.AreaLaw.LocalSolution.exists_unit_mem_coordinateSliceProjection_of_box,
        TNLean.MPS.AreaLaw.LocalSolution.card_le_two_mul_re_trace_of_coordinateSliceProjection_of_box}
    \leanok
    \uses{thm:hastings05_coordinate_slice_approximation}
    Let $\Lambda=\prod_{k=1}^D\{0,\ldots,L_k-1\}$, where $D\ge1$ and
    every $L_k\ge1$. Choose any coordinate axis $a$ and take $B=L_a$.
    Then $B\le|\Lambda|$, the slices of width $2\ell$ cover the box, and
    maximum-metric locality implies strip locality. Hence all conclusions
    of the preceding theorem hold, including the cubic logarithmic length
    bound in the physical number of sites, in every dimension.
\end{corollary}
\begin{proof}\leanok
    Vary coordinate $a$ while fixing the other coordinates at zero to
    inject its side into the full box. Also $z(v)<L_a\le2\ell L_a$.
    The coordinate projection decreases maximum-metric distance.
\end{proof}
